\section{Character Creation}

In \game{} you play a party of three characters. Each character is created based on an \kw{archetype}.
All characters also share the Basic Actions listed in Combat. The three archetypes are:

\begin{itemize}
    \item \textbf{Warrior:} Front-line steel and grit. Holds ground when the Shard-dark presses in.
    \begin{itemize}
        \item \kw{Cleave} --- \hard{} \hard{}: Deal 1 hit to up to two foes in your zone.
        \item \kw{Hold the Line} --- \medium{} \energy{1}: You and one ally in your zone each \kw{Guard} (ignore the first hit).
    \end{itemize}
    \item \textbf{Rogue:} Fast blades and cruel timing. Strikes, then vanishes into another zone.
    \begin{itemize}
        \item \kw{Skirmish} --- \medium{} \energy{1}: Deal 1 hit to a foe in your zone, then \kw{Move} to an adjacent zone.
        \item \kw{Exploit} --- \hard{} \hard{}: Deal 2 hits to one foe in your zone that was already hit this round.
    \end{itemize}
    \item \textbf{Mage:} Shard-touched will. Channels Aether-Glass into bolts and wards.
    \begin{itemize}
        \item \kw{Arc Bolt} --- \medium{} \energy{1}: Deal 1 hit to a foe in your zone or an adjacent zone.
        \item \kw{Ward} --- \easy: One ally in your zone gains \kw{Guard} (ignore the first hit).
    \end{itemize}
\end{itemize}

Then choose one \kw{class} under your archetype (see Classes). You begin at level 1: pick one of that level's five special actions. Each time you level up, pick one special action from the new level's list of five.

\subsection{Dice Pool and Level}

Every character starts with a combat \kw{Dice Pool} of \D{2}{6}. The pool grows by one die at \textbf{level 3}, \textbf{level 6}, and \textbf{level 9} (\D{3}{6}, \D{4}{6}, \D{5}{6}). A bigger pool means more dice to assign in a round and, from level 3 on, access to \textbf{power actions} that need three or four dice from one character: Grand \grand{} at level 3+, Dire \dire{} at level 6+, Fell \fell{} at level 9+. The class lists only offer those actions at the levels that can use them.

\subsection{Starting Kit}

Each character begins with one \textbf{weapon} of their choice, \textbf{Light armor}, and \textbf{one gear item} from the Equipment chapter. One character may take a Shield instead of the gear item. A Mage also begins with a \textbf{Grimoire} holding two First-Circle spells (see Spells). Everything else is bought with Refined crystals in a Town or taken from the dark.

\subsection{Attributes}

Each character has six attributes. In combat they do nothing---the Dice Pool carries the fight. Everywhere else, they are who your character is when the dice stop.

\begin{itemize}
    \item \textbf{Strength $\left[\textbf{STR}\right]$:} Hauling sleds, fording black water, forcing doors, holding a beam while the others crawl under.
    \item \textbf{Dexterity $\left[\textbf{DEX}\right]$:} Sneaking past tolls, picking Guild locks, scouting ahead, mending straps and tools by lamplight.
    \item \textbf{Constitution $\left[\textbf{CON}\right]$:} Marching on a filled clock, breathing ash, resisting plague banners. CON sets your \kw{Hits}.
    \item \textbf{Intelligence $\left[\textbf{INT}\right]$:} Reading pre-Cataclysm script, forging papers, recalling what a Gate did to the last expedition.
    \item \textbf{Wisdom $\left[\textbf{WIS}\right]$:} Noticing the rotated milestone, reading Omens in coals, tending wounds without making them worse.
    \item \textbf{Charisma $\left[\textbf{CHA}\right]$:} Bargaining with factors, bribing ferrymen, talking a collector into one more week.
\end{itemize}

\subsection{Attribute Dice}

Each attribute is a number of \kw{Attribute Dice} from 1 to 4---that is the whole score.
Roll \D{3}{6} once for each of the six attributes and read the dice from the table.
No attribute may ever have more than 4 dice.

\begin{table}[htbp]
  \centering
  \caption{Determine Attribute Dice}
  \label{tab:attribute_dice}
\renewcommand{\arraystretch}{1.25}
\begin{tabular}{cc}
  \shardheadercell{3d6} & \shardheadercell{Dice} \\
  \rowcolor{bone} 3 -- 4   & 1 \\
  \rowcolor{ash!15} 5 -- 8   & 2 \\
  \rowcolor{bone} 9 -- 16  & 3 \\
  \rowcolor{ash!15} 17 -- 18 & 4 \\
\end{tabular}
\end{table}

After rolling, you may swap the dice of any two attributes once, so the Warrior is not the party's only scholar by accident.

Each time a character \kw{level up}s, add $+1$ die to one attribute (maximum 4).
Write the dice on the party sheet; they are what you roll on the road.

\subsection{Hits}

A character's maximum \kw{Hits} are 3 plus their CON Attribute Dice.
CON 1 die gives 4 Hits; 2 dice give 5; 3 dice give 6; 4 dice give 7.
Each \kw{Scar} earned by going Down in a fight lowers the maximum by 1, never below 3.
Hits are tracked on the party sheet and heal only through Campfire Counsel, class actions, or encounter rewards.

\subsection{Attribute Check}

Outside combat the game asks for \kw{Attribute Check}s. One character attempts the check: roll their \kw{Attribute Dice} for the fitting attribute and count every die showing \textbf{4 or more} as a \kw{success}.

Encounters, Omens, and Campfire Counsel name their difficulty with the same words as combat. Outside combat those words are the successes you need:

\begin{center}
{\small
\renewcommand{\arraystretch}{1.2}
\begin{tabularx}{\linewidth}{l c >{\raggedright\arraybackslash}X}
  \shardheadercell{Difficulty} & \shardheadercell{Successes} & \shardheadercell{Example} \\
  \rowcolor{bone} Easy & 1 & Notice a leaning milestone \\
  \rowcolor{ash!15} Medium & 2 & Sneak past a mercenary toll \\
  \rowcolor{bone} Hard (``Hard Hard'') & 3 & Forge quarantine papers \\
\end{tabularx}
}
\end{center}

\begin{itemize}
    \item \textbf{Easy floor:} an Easy check always rolls at least 2 dice, even for an attribute with only 1 die.
    \item \kw{Aid}\textbf{:} before the roll, each other standing character may spend \energy{1} from the party track to add 1 die to the pool. Each may Aid once; their own attributes do not matter.
    \item \textbf{Energy in the text:} a check written ``Medium + Energy'' costs that Energy on top of any Aid, spent before rolling.
    \item \kw{Exhausted} removes 1 die from the pool; a \kw{Down} character can neither attempt nor Aid.
\end{itemize}

\paragraph{Which attribute?}
If the text names an attribute, use it. If it only describes the task, pick from this table; if two fit, use the higher.

\begin{center}
{\small
\renewcommand{\arraystretch}{1.15}
\begin{tabularx}{\linewidth}{l >{\raggedright\arraybackslash}X}
  \shardheadercell{Attribute} & \shardheadercell{Tasks} \\
  \rowcolor{bone} STR & climb, ford, swim, haul, hold a line, force \\
  \rowcolor{ash!15} DEX & sneak, scout, pick locks, repair, duel, raft-build \\
  \rowcolor{bone} CON & march, endure heat or cold, resist plague or fumes, camp cold \\
  \rowcolor{ash!15} INT & lore, maps, forge papers, read script, appraise glass \\
  \rowcolor{bone} WIS & notice, insight, heal, read omens or coals, keep watch \\
  \rowcolor{ash!15} CHA & bargain, bribe, negotiate passage, sing, lead refugees \\
\end{tabularx}
}
\end{center}

\begin{rpgbox}{Attribute Check Example}
    \textit{Medium crossing (STR):} your Warrior has \attr{STR}{3} dice. You roll $5$, $4$, $2$---two successes. The party crosses.

    \textit{Easy notice (WIS):} your Mage has \attr{WIS}{2} dice. You roll $3$, $1$---no success. The milestone was rotated, and you walk the wrong way.

    \textit{Hard Hard forgery (INT):} the same Mage has \attr{INT}{3} dice; alone that is a 1-in-8 chance of three successes. The Warrior and Rogue each spend \energy{1} to Aid, making 5 dice. You roll $6$, $6$, $4$, $3$, $1$---three successes, papers accepted, Energy down by 2.
\end{rpgbox}

\subsubsection{Critical Success}
If you rolled at least 2 dice and \textbf{every} die shows a 6, you score a \kw{Critical Success}. Apply the encounter's success and also \textbf{clear 1 Resonance Stress}, unless the text grants a better bonus.

\subsubsection{Critical Miss}
If you rolled at least 2 dice and \textbf{every} die shows a 1, it is a \kw{Critical Miss}. Apply the failure and also \textbf{mark $+1$ Resonance Stress}, unless the text names a worse cost.

\paragraph{Spending Favor.}
After rolling, you may burn one Favor mark with any faction to reroll every die that did not show a success; keep the new results. One Favor per check.

\subsection{Party}

Your three characters live on one party sheet (see the last page of this book). Record:

\begin{itemize}
    \item \textbf{Each character:} name, archetype, class, level, Dice Pool, six attributes (1--4 dice each), maximum and current Hits, Scars, archetype actions, the one special action learned per level, weapon, off-hand, armor and its boxes, four gear slots, bound relic, and (Mages) known spells.
    \item \textbf{Energy} (8 boxes), shared by all three.
    \item \textbf{Resonance Stress} (6 boxes) --- see Non-Combat Encounters and Campaign: Corruption Mutations.
    \item \textbf{Debt and Favor} marks, one line per faction (Guild, Crown, Cult, Companies, others as they appear).
    \item \textbf{Journey Clock} for the current road and the \textbf{World Crack Clock} (8 boxes).
\end{itemize}

\paragraph{Bonds.}
Name one tie between each pair of characters: a debt, a shared grave, a lie one told the other. Bonds have no rule weight; they are what you read back when a Scar is written down, and they decide who carries whom out of the Gate.
